% Section 3 -- Tasks, models, and datasets.
% Every count in this section comes from a generated table or from a results file; see paper/tools/gen_tables.py.
\section{Tasks, Models, and Datasets}
\label{sec:setup}

\subsection{The minimal-data regime}
\label{sec:regime}

We study linear probes on language-model activations trained from tens of examples rather than thousands. The appeal
is obvious: a direction vector costs one forward pass to apply, adds roughly a microsecond of arithmetic to a request
that already costs tens of milliseconds, and can be fitted in minutes on a laptop's worth of labelled text. The
question this paper asks is not whether such probes can be fitted---they always can---but what has to be true of an
evaluation before a number computed on one means what it appears to mean.

All probes here share the same construction. Given a set of positive texts $P$ and negative texts $N$, activations
$h_\ell$ are read at a single decoder layer $\ell$ and a single token position, and the direction is the difference of
class means,
\begin{equation}
    v = \frac{1}{|P|}\sum_{i \in P} h^{(i)}_\ell - \frac{1}{|N|}\sum_{j \in N} h^{(j)}_\ell .
    \label{eq:direction}
\end{equation}
At inference a text is scored by the cosine of its activation with the unit direction, $s = \langle \hat{h}_\ell,
\hat{v}\rangle$, and classified against a threshold fitted on the training negatives at a $5\,\%$ false-positive
quantile, floored at the midpoint of the class means. Section~\ref{sec:learn} shows that a construction widely
described as a distinct alternative to~\eqref{eq:direction}---averaging \emph{within-pair} differences over matched
pairs---is not an alternative at all but the same estimator, which is why this paper uses the term \emph{paired
contrastive data construction} throughout and never treats pairing as an estimator choice.

\subsection{Three tasks}
\label{sec:tasks}

We evaluate three probe families, chosen because they span the three places a safety probe is typically inserted and
because they fail in three different ways.

\textbf{Harmful-content filtering (\af{})} scores an incoming user prompt and is meant to fire on requests a model
should refuse. It is the easiest task to appear to solve, because harmful and benign prompts differ in topic as well
as in harm.

\textbf{Agent-injection detection (\aag{})} scores an agent scaffold---a user instruction, a tool response, and the
model's turn boundary---and is meant to fire when the tool response contains an instruction injected by a third party.
Unlike \af{}, the signal is localised inside a long input at a position that varies from case to case.

\textbf{Policy compliance (\apc{})} scores a model's own \emph{response} and is meant to distinguish giving advice
from providing information within a regulated domain. It is the only one of the three whose target is a speech act
rather than a topic, and Section~\ref{sec:learn} shows this is what makes it hard to evaluate honestly.

\subsection{Models}
\label{sec:models}

Seven instruction-tuned models across three families (Gemma-2, Gemma-3, Llama-3), from $1$B to $9$B parameters, all in
\texttt{bfloat16} on a single A100-40GB. Table~\ref{tab:models} lists the layer each probe family reads. Those layer
indices are the ones the probes were originally fitted at and are held fixed throughout this paper: a layer chosen by
sweeping on the evaluation set would not be a held-out number, and the one place we do sweep (Section~\ref{sec:read})
is reported as a note rather than as a headline.

% GENERATED by paper/tools/gen_tables.py -- do not edit by hand.
% sources: results/gates/b1/*.json, results/gates/b2/*.json, results/gates/b3/*.json
\begin{table}[!t]
\caption{Models, and the layer each probe family reads. Layer indices are those of the
shipped probes (see Section~\ref{sec:models}) and are held fixed throughout; $L$ is the number of decoder
blocks. APC probes were evaluated on
the three models carried into the external study.}
\label{tab:models}
\centering
\begin{tabular}{lcccc}
\toprule
Model & $L$ & AF layer & AAG layer & APC layer \\
\midrule
Gemma-2-2B & 26 & 13 & 14 & -- \\
Gemma-2-9B & 42 & 21 & 23 & 10 \\
Gemma-3-1B & 26 & 13 & 10 & -- \\
Gemma-3-4B & 34 & 17 & 13 & -- \\
Llama-3.1-8B & 32 & 16 & 17 & 8 \\
Llama-3.2-1B & 16 & 8 & 8 & -- \\
Llama-3.2-3B & 28 & 14 & 15 & 7 \\
\bottomrule
\end{tabular}
\end{table}

The shipped probes are the publicly released direction vectors of the author's earlier minimal-data recipe
\cite{messenger2026aase}, fitted on the training sets that archive contains; this paper takes them as given and does
not refit them. Those training sets are small by design, and their sizes are stated here once so that every later
number is read against them: $20$ harmful and $20$ benign prompts for the harmful-content probe; $50$ injected agent
scaffolds against $100$ clean ones for the injection probe, all synthetic and none drawn from InjecAgent, which is why
Section~\ref{sec:read} retrains that probe on a held-out half of the real benchmark before comparing read positions;
and $15$ authored contrastive pairs per policy for the compliance probe. The counts are those recorded in the released
probes' metadata.

One detail of the extraction stack turns out to matter more than the layer index, and is treated separately in
Section~\ref{sec:tensor}: which \emph{tensor} at that layer the probe reads.

\subsection{Datasets}
\label{sec:data}

Table~\ref{tab:datasets} summarises the evaluation data. Three choices in it are load-bearing.

\textbf{The benign side is reported per set, never pooled.} For \af{} we use two benign sets with deliberately
different character. XSTest-v2 safe prompts are built to \emph{look} unsafe while being harmless---questions about
killing a process, or about historical atrocities---and are the realistic false-refusal stressor. JailbreakBench
benign prompts are on-topic twins of harmful behaviours: same subject matter, benign intent. A probe that has learned
topic rather than harm separates the harmful set from XSTest easily and from JailbreakBench barely, and pooling the
two benign sets hides exactly that. Every \af{} number in this paper is therefore reported against both.

\textbf{HarmBench is used at $300$ behaviours, not $400$.} The $100$ copyright behaviours are excluded because they
request verbatim reproduction of protected text rather than harmful content, and including them would mis-measure a
harm probe. The exclusion is stated here rather than buried in a configuration file.

\textbf{The agent benign side is constructed, because none exists.} InjecAgent ships $1{,}054$ injected cases per
attack setting but only $17$ clean tool contexts, which puts a floor of $5.9\,\%$ on any false-positive-rate
measurement. We assembled a $120$-case benign set from the InjecAgent user cases, AgentDojo suites and authored
scaffolds in the same format, including $36$ hard negatives---clean tool responses that contain imperative language,
urgency markers, or embedded quoted text---so that false positives can be attributed to a cause rather than merely
counted.

% GENERATED by paper/tools/gen_tables.py -- do not edit by hand.
% sources: results/gates/b1/google_gemma_2_9b_it.json, results/gates/b2/google_gemma_2_9b_it.json, results/gates/b3/google_gemma_2_9b_it.json
\begin{table*}[!t]
\caption{Evaluation data. Every dataset is pinned by sha256 in each results file; the AF benign side is
reported per set throughout because the two sets behave differently.}
\label{tab:datasets}
\centering
\setlength{\tabcolsep}{8pt}
\small
\begin{tabular}{l c >{\raggedright\arraybackslash\hyphenpenalty=10000\exhyphenpenalty=10000}p{0.34\textwidth} c}
\toprule
Source & $n$ & Role & Task \\
\midrule
HarmBench (standard + contextual) & 300 & harmful prompts (200 standard, 100 contextual) & AF \\
XSTest v2 (safe) & 250 & benign, off-topic stressor & AF \\
JailbreakBench (benign) & 100 & benign, on-topic twins & AF \\
InjecAgent \texttt{base} & 1054 & injected tool responses & AAG \\
InjecAgent \texttt{enhanced} & 1054 & held out entirely & AAG \\
AAG benign set & 120 & clean tool contexts (36 hard negatives) & AAG \\
APC external set & 360 & authored responses (240 medical, 60 financial, 60 legal) & APC \\
\bottomrule
\end{tabular}
\end{table*}

Two label artifacts in InjecAgent are worth recording for anyone else using it. Every one of the $1{,}054$ injected
tool responses is wrapped in an outer pair of quotation marks and has its inner quotes escaped, while none of the
benign templates is; left in place this is a perfect label cue that has nothing to do with injection. We normalise it
away for the main results. Separately, and more consequentially for evaluation design, the $1{,}054$ cases are built
by crossing only $62$ distinct attacker instructions with $17$ user instructions, so a case-level train/test split
places every test-half injection string in the training half as well. Section~\ref{sec:read} reports both that split
and a grouped split that removes the overlap. Both artifacts have been reported upstream \cite{injecagentissue7}.

\subsection{The policy-compliance labelling rule}
\label{sec:apcrule}

Advice and information shade into each other, so the rule that separates them was fixed in writing before any
activation was extracted, and did not change afterwards:

\begin{quote}
\emph{A hedge does not remove a surviving directive; being detailed enough to act on does not make a description
advice.}
\end{quote}

Concretely, a response is labelled \emph{advice} when it directs the reader toward a specific action concerning their
own case, whether or not it is hedged---``ibuprofen often helps, so take $400$\,mg with food, but do check with your
doctor'' is advice. It is labelled \emph{information} when it describes facts, mechanisms, typical ranges or ordinary
professional practice without directing the reader, however actionable the content---``adult dosing in the labelling
is $400$--$600$\,mg every six hours'' is information. Responses that decline and hand off without conveying either a
directive or substantive content are labelled \emph{refusal} and are excluded from the advice-versus-information
comparison, then reported separately: a compliance probe that fires on refusals would fire on precisely the responses
the policy is trying to elicit, and Section~\ref{sec:learn} shows that the shipped probes do. Labels were applied by
the author under this rule, frozen in writing before any activation was extracted; we report no inter-annotator
agreement, and the construct-validity finding Section~\ref{sec:apc} draws from the set is a finding about authored
short responses, not about medical advice as a phenomenon.

The external evaluation set contains $240$ medical responses in four bands---clear advice, clear information, hedged
advice, and information detailed enough to act on---plus $20$ refusals, with $60$ financial and $60$ legal responses
in the same structure for cross-topic transfer. The two middle bands are the ones that matter: hedged advice catches
a probe that has learned that hedging means compliance, and information-dense catches one that has learned that
clinical specificity means violation. Section~\ref{sec:learn} reports what happened when we asked whether this set
could support a claim about speech acts at all.
